\documentclass[11pt]{article}

\usepackage[margin=0.62in]{geometry}
\usepackage{amsmath,amssymb}
\usepackage{enumitem}
\usepackage{xcolor}
\usepackage{tcolorbox}
\usepackage{fancyhdr}
\usepackage{lastpage}
\usepackage{hyperref}
\usepackage{microtype}
\usepackage{array}

\definecolor{uwpurple}{HTML}{4B2E83}
\definecolor{uwgold}{HTML}{B7A57A}
\definecolor{softpurple}{HTML}{F4F1FA}
\definecolor{softgold}{HTML}{FBF7EA}
\definecolor{deepgold}{HTML}{8A6F2A}

\tcbuselibrary{breakable,skins}

\hypersetup{
    colorlinks=true,
    linkcolor=uwpurple,
    urlcolor=uwpurple
}

\pagestyle{fancy}
\fancyhf{}
\lhead{\textbf{MATH 394: Probability I}}
\rhead{\textbf{Final Examination}}
\cfoot{\thepage\ of \pageref{LastPage}}
\setlength{\headheight}{14pt}

\setlength{\parindent}{0pt}
\setlength{\parskip}{0.4em}

\newcommand{\course}{MATH 394: Probability I}
\newcommand{\instructor}{Arman Jahangiri}
\newcommand{\department}{Department of Mathematics}
\newcommand{\university}{University of Washington}

\newtcolorbox{problemheader}{
    colback=softpurple,
    colframe=uwpurple,
    boxrule=0.8pt,
    arc=2mm,
    left=2mm,
    right=2mm,
    top=1mm,
    bottom=1mm
}

\newtcolorbox{policybox}{
    breakable,
    colback=softgold,
    colframe=deepgold,
    boxrule=0.8pt,
    arc=2mm,
    left=2.5mm,
    right=2.5mm,
    top=1.5mm,
    bottom=1.5mm
}

\newcommand{\workspace}[1]{%
    \vspace{0.4cm}
    \begin{tcolorbox}[
        height=#1,
        colback=white,
        colframe=uwgold,
        boxrule=0.5pt,
        arc=1mm,
        title={\small\textbf{Work space}},
        coltitle=uwpurple,
        colbacktitle=softpurple
    ]
    \end{tcolorbox}
}

\begin{document}

\begin{titlepage}
\centering
\vspace*{0.05cm}

{\Large \textbf{\university}}\\
{\large \department}

\vspace{0.28cm}

{\LARGE \color{uwpurple}\textbf{Final Examination}}\\[0.15cm]
{\Large \textbf{\course}}\\[0.15cm]
{\large Summer 2026 \quad | \quad Instructor: \instructor}

\vspace{0.22cm}

\begin{tcolorbox}[
    width=0.94\textwidth,
    colback=softpurple,
    colframe=uwpurple,
    boxrule=0.9pt,
    arc=3mm
]
\centering
\renewcommand{\arraystretch}{1.15}
\begin{tabular}{>{\bfseries}p{0.19\textwidth}p{0.68\textwidth}}
Date: & August 21, 2026 \\
Answering time: & 60 minutes \\
Total points: & 60 \\
\end{tabular}
\end{tcolorbox}

\vspace{0.18cm}

\begin{tcolorbox}[
    width=0.94\textwidth,
    colback=white,
    colframe=uwgold,
    boxrule=0.8pt,
    arc=3mm
]
\renewcommand{\arraystretch}{1.35}
\begin{tabular}{>{\bfseries}p{0.22\textwidth}p{0.67\textwidth}}
Name: & \hrulefill \\
Student ID: & \hrulefill \\
Signature: & \hrulefill \\
\end{tabular}
\end{tcolorbox}

\vspace{0.15cm}

\begin{policybox}
\small
\textbf{Instructions}
\begin{itemize}[leftmargin=*,itemsep=0.05em,topsep=0.15em,parsep=0pt]
    \item This examination contains \textbf{seven problems}. Verify that your copy is complete.
    \item Show all work and justify your answers. Unsupported answers may receive little or no credit.
    \item Clearly define any events or random variables that you introduce.
    \item The estimated times are provided only as a guide for pacing.
    \item Do not remove pages from the examination booklet.
\end{itemize}

\medskip
\textbf{Permitted materials}

Examinations are closed-book and closed-notes. Students may use:
\begin{itemize}[leftmargin=*,itemsep=0.05em,topsep=0.12em,parsep=0pt]
    \item one handwritten 8.5 $\times$ 11 inch reference sheet (front and back), and
    \item a non-programmable non-graphing calculator capable of standard exponential and logarithmic computations (e.g. TI-30X IIS).
\end{itemize}
No other electronic devices, calculators, communication devices, or external resources may be used during examinations.
\end{policybox}

\vspace{0.10cm}

\renewcommand{\arraystretch}{1.18}
\begin{tabular}{|>{\centering\arraybackslash}p{0.16\textwidth}|>{\centering\arraybackslash}p{0.16\textwidth}|}
\hline
\textbf{Problem} & \textbf{Score} \\
\hline
1 & /16 \\
\hline
2 & /8 \\
\hline
3 & /7 \\
\hline
4 & /7 \\
\hline
5 & /8 \\
\hline
6 & /8 \\
\hline
7 & /6 \\
\hline
\textbf{Total} & \textbf{/60} \\
\hline
\end{tabular}

\vspace{0.08cm}
{\footnotesize By signing above, you affirm that the work is your own and that you followed the examination rules.}
\end{titlepage}



% ------------------------------------------------------------
\begin{problemheader}
\textbf{Problem 1.} 
\hfill
\textbf{Estimated time: 13 min}
\qquad
\textbf{16 points}
\end{problemheader}

Let the joint probability density function of random variables \(X\) and \(Y\)
be given by
\[
f(x,y)
=
\begin{cases}
x^2e^{-x(y+1)}, & x\geq0,\quad y\geq0,\\
0, & \text{elsewhere}.
\end{cases}
\]

The marginal probability density functions are
\[
f_X(x)
=
\begin{cases}
xe^{-x}, & x\geq0,\\
0, & \text{otherwise},
\end{cases}
\]
and
\[
f_Y(y)
=
\begin{cases}
\dfrac{2}{(1+y)^3}, & y\geq0,\\
0, & \text{otherwise}.
\end{cases}
\]

\begin{enumerate}[label=(\alph*),leftmargin=*,itemsep=0.8em]

    \item Calculate
    \[
    \operatorname{Cov}(X,Y).
    \]
    \hfill \textbf{(4 points)}

    \item Calculate
    \[
    \operatorname{Cov}(3X+1,\,2X-4).
    \]
    \hfill \textbf{(3 points)}

    \item Are \(X\) and \(Y\) independent? Justify your answer.
    \hfill \textbf{(2 points)}

    \item Calculate the moment-generating function of \(X\), namely
    \[
    M_X(t)=E(e^{tX}).
    \]
    \hfill \textbf{(5 points)}

    \item Using \(M_X(t)\), calculate \(E(X)\).
    \hfill \textbf{(2 points)}

\end{enumerate}

\hrule



\newpage

% ------------------------------------------------------------
\begin{problemheader}
\textbf{Problem 2.}
\hfill
\textbf{Estimated time: 9 min}
\qquad
\textbf{8 points}
\end{problemheader}

Let \(X\) and \(Y\) be two positive independent continuous random variables
with probability density functions \(f_1(x)\) and \(f_2(y)\), respectively.

Find the probability density function of
\[
U=\frac{X}{Y}.
\]

\textit{Hint:} Let \(V=X\). Find the joint probability density function of
\(U\) and \(V\) (Jacobian method), and then calculate the marginal probability
density function of \(U\).

\hrule



\newpage

% ------------------------------------------------------------
\begin{problemheader}
\textbf{Problem 3.}
\hfill
\textbf{Estimated time: 5 min}
\qquad
\textbf{7 points}
\end{problemheader}

Let \(X\) and \(Y\) be i.i.d.\ positive random variables. Assume that all expectations appearing below exist and are finite.

For each part below, fill in the appropriate equality or inequality symbol.

Write \(=\) if the two sides are always equal, \(\leq\) if the left-hand side
is less than or equal to the right-hand side but not necessarily equal, and
similarly for \(\geq\). If no relation holds in general, write \(?\).

Justify each answer.

\begin{enumerate}[label=(\alph*),leftmargin=*,itemsep=1.2em]

    \item
    \[
    E(\log X)
    \quad \underline{\hspace{1.2cm}} \quad
    \log(E(X)).
    \]
    \hfill \textbf{(3 points)}

    \item
    \[
    P(X\leq Y)
    \quad \underline{\hspace{1.2cm}} \quad
    P(X\geq Y).
    \]
    \hfill \textbf{(2 points)}

    \item
    \[
    E(XY)
    \quad \underline{\hspace{1.2cm}} \quad
    \sqrt{E(X^2)E(Y^2)}.
    \]
    \hfill \textbf{(2 points)}

\end{enumerate}

\hrule



\newpage

% ------------------------------------------------------------
\begin{problemheader}
\textbf{Problem 4.}
\hfill
\textbf{Estimated time: 8 min}
\qquad
\textbf{7 points}
\end{problemheader}

Suppose
\[
X_1,\ldots,X_n
\]
are independent Bernoulli random variables with unknown success probability
\(p\), and define
\[
\widehat p
=
\frac1n\sum_{i=1}^n X_i.
\]

Let
\[
\varepsilon>0
\qquad\text{and}\qquad
0<\alpha<1.
\]

Using Chebyshev's inequality, derive a condition on the sample size \(n\) that
guarantees
\[
P\left(|\widehat p-p|<\varepsilon\right)
\geq
1-\alpha
\]
for every
\[
0<p<1.
\]

\hrule



\newpage

% ------------------------------------------------------------
\begin{problemheader}
\textbf{Problem 5.}
\hfill
\textbf{Estimated time: 7 min}
\qquad
\textbf{8 points}
\end{problemheader}

Let \(X\) and \(Y\) be independent exponential random variables with common
parameter \(\lambda>0\):
\[
X,Y\overset{\mathrm{i.i.d.}}{\sim}\operatorname{Exp}(\lambda),
\]
with probability density function
\[
f(w)=
\begin{cases}
\lambda e^{-\lambda w}, & w>0,\\
0, & \text{otherwise}.
\end{cases}
\]

Define
\[
T=X+Y.
\]

Using the convolution formula
\[
f_T(t)
=
\int_{-\infty}^{\infty}
f_X(x)f_Y(t-x)\,dx,
\]
find the probability density function of \(T\), including its support.

\hrule



\newpage

% ------------------------------------------------------------
\begin{problemheader}
\textbf{Problem 6.}
\hfill
\textbf{Estimated time: 5 min}
\qquad
\textbf{8 points}
\end{problemheader}

Let \(X,Y,Z\) be jointly continuous with joint probability density function
\[
f_{X,Y,Z}(x,y,z)
=
\begin{cases}
x^2e^{-x(1+y+z)}, & x,y,z>0,\\
0, & \text{otherwise}.
\end{cases}
\]

The marginal probability density functions are
\[
f_X(x)
=
\begin{cases}
e^{-x}, & x>0,\\
0, & \text{otherwise},
\end{cases}
\]
and
\[
f_Y(y)
=
\begin{cases}
\dfrac{1}{(1+y)^2}, & y>0,\\
0, & \text{otherwise},
\end{cases}
\qquad
f_Z(z)
=
\begin{cases}
\dfrac{1}{(1+z)^2}, & z>0,\\
0, & \text{otherwise}.
\end{cases}
\]

\begin{enumerate}[label=(\alph*),leftmargin=*,itemsep=0.8em]

    \item Are \(X,Y,Z\) pairwise independent? Justify your answer.
    \hfill \textbf{(5 points)}

    \item Are \(X,Y,Z\) mutually independent? Justify your answer.
    \hfill \textbf{(3 points)}

\end{enumerate}

\hrule



\newpage

% ------------------------------------------------------------
\begin{problemheader}
\textbf{Problem 7.}
\hfill
\textbf{Estimated time: 4 min}
\qquad
\textbf{6 points}
\end{problemheader}

For each part below, select the \textbf{correct statement} of the indicated
theorem. Only one choice is correct in each part.

\begin{enumerate}[label=(\alph*),leftmargin=*,itemsep=1.2em]

% ------------------------------------------------------------
\item \textbf{Weak Law of Large Numbers (WLLN).}

Let \(X_1,X_2,\ldots\) be i.i.d.\ random variables with finite mean
\[
E(X_i)=\mu
\]
and finite variance
\[
\operatorname{Var}(X_i)=\sigma^2>0.
\]
Let
\[
\overline X_n=\frac1n\sum_{i=1}^n X_i.
\]

Which of the following is the conclusion of the Weak Law of Large Numbers?
\hfill \textbf{(2 points)}

\begin{enumerate}[label=(\Alph*),leftmargin=2em,itemsep=0.5em]

    \item For every \(\varepsilon>0\),
    \[
    P\left(
    |\overline X_n-\mu|>\varepsilon
    \right)=0
    \qquad
    \text{for every }n.
    \]

    \item For every \(\varepsilon>0\),
    \[
    \lim_{n\to\infty}
    P\left(
    |\overline X_n-\mu|>\varepsilon
    \right)
    =0.
    \]

    \item
    \[
    P\left(
    \lim_{n\to\infty}\overline X_n=\mu
    \right)=1.
    \]

    \item For every \(z\in\mathbb R\),
    \[
    \lim_{n\to\infty}
    P\left(
    \frac{\overline X_n-\mu}{\sigma/\sqrt n}
    \leq z
    \right)
    =
    \Phi(z).
    \]

\end{enumerate}


% ------------------------------------------------------------
\item \textbf{Strong Law of Large Numbers (SLLN).}

Let \(X_1,X_2,\ldots\) be i.i.d.\ random variables satisfying the same
assumptions as in the Weak Law of Large Numbers, with
\[
E(X_i)=\mu
\]
and finite variance
\[
\operatorname{Var}(X_i)=\sigma^2>0.
\]

Which of the following is the conclusion of the Strong Law of Large Numbers?
\hfill \textbf{(2 points)}

\begin{enumerate}[label=(\Alph*),leftmargin=2em,itemsep=0.5em]

    \item
    \[
    P\left(
    \lim_{n\to\infty}\overline X_n=\mu
    \right)=1.
    \]

    \item For every \(\varepsilon>0\),
    \[
    \lim_{n\to\infty}
    P\left(
    |\overline X_n-\mu|>\varepsilon
    \right)
    =0.
    \]

    \item
    \[
    \lim_{n\to\infty}
    P(\overline X_n=\mu)
    =1.
    \]

    \item
    \[
    \frac{\overline X_n-\mu}{\sigma/\sqrt n}
    \xrightarrow[]{d}
    N(0,1).
    \]

\end{enumerate}


% ------------------------------------------------------------
\item \textbf{Central Limit Theorem (CLT).}

Let \(X_1,X_2,\ldots\) be i.i.d.\ random variables with
\[
E(X_i)=\mu,
\qquad
\operatorname{Var}(X_i)=0<\sigma^2<\infty.
\]

Which of the following is the conclusion of the Central Limit Theorem?
\hfill \textbf{(2 points)}

\begin{enumerate}[label=(\Alph*),leftmargin=2em,itemsep=0.5em]

    \item
    \[
    \overline X_n
    \xrightarrow[]{P}
    N(\mu,\sigma^2).
    \]

    \item
    \[
    P\left(
    \lim_{n\to\infty}\overline X_n=\mu
    \right)=1.
    \]

    \item For every \(z\in\mathbb R\),
    \[
    \lim_{n\to\infty}
    P\left(
    \frac{\overline X_n-\mu}{\sigma/\sqrt n}
    \leq z
    \right)
    =
    \Phi(z).
    \]

    \item For every \(\varepsilon>0\),
    \[
    \lim_{n\to\infty}
    P\left(
    |\overline X_n-\mu|>\varepsilon
    \right)
    =0.
    \]

\end{enumerate}

\end{enumerate}

\hrule



\end{document}